\section{Introduction}

Every decade following the U.S. Census, states must redraw congressional district boundaries to reflect population changes. This redistricting process determines electoral representation for 330 million Americans across 435 congressional seats, yet it remains one of the most contentious aspects of American democracy \cite{daley2021unrigging}. Despite constitutional requirements for equal population and Supreme Court mandates for fairness \cite{reynolds1964,shaw1993}, the practice of partisan gerrymandering—manipulating district boundaries for electoral advantage—undermines public confidence in representative government.

The 2019 Supreme Court decision in \textit{Rucho v. Common Cause} \cite{rucho2019} declared partisan gerrymandering claims non-justiciable, effectively placing redistricting beyond federal judicial oversight. This ruling heightened calls for alternative approaches that remove human discretion from boundary drawing. While independent redistricting commissions offer one solution, they remain vulnerable to political pressure and lack transparency in decision-making criteria \cite{stephanopoulos2015partisan}.

Mathematical and computational approaches to redistricting are not new \cite{hess1965nonpartisan,garfinkel1970optimal}. However, recent advances in graph partitioning algorithms and computational power enable large-scale implementation across entire states. More importantly, these approaches can embody the same principles of mathematical fairness that have successfully governed congressional apportionment for over 80 years.

In 1941, Congress adopted the Huntington-Hill method \cite{huntington1928method} to apportion the 435 House seats among states. This method, based on minimizing relative differences in representation, has operated without controversy for eight decades precisely because it is algorithmic, objective, and mathematically defensible \cite{balinski1982fair}. No state can credibly claim unfair treatment when seat allocation follows from transparent formulae applied uniformly nationwide.

We argue that similar principles can govern district boundary design. If mathematical objectivity resolved the politically charged question of \textit{how many} seats each state receives, why not apply the same philosophy to \textit{where} district boundaries should be drawn? This paper presents an algorithmic redistricting method based on recursive graph bisection that extends Huntington-Hill's commitment to mathematical rigor from apportionment to boundary design.

\subsection{Our Approach}

We model each state as a graph where census tracts form nodes and geographic adjacency defines edges. Using the METIS graph partitioning library \cite{karypis1998fast}, we recursively bisect this graph to create congressional districts that:
\begin{itemize}
    \item Achieve near-perfect population equality (mean deviation: 2.79\%)
    \item Maintain geographic contiguity by construction
    \item Optimize compactness through edge-cut minimization
    \item Produce reproducible results from the same input data
    \item Contain zero partisan intent by design
\end{itemize}

We apply this method to all 50 states using 2020 Census data \cite{census2020redistricting} and 2020 presidential election results \cite{ansolabehere2012mecubed}, analyzing the resulting 435 districts for population balance, compactness, political characteristics, and demographic representation. Results demonstrate that algorithmic redistricting can simultaneously satisfy traditional criteria—population equality, contiguity, and compactness—while maintaining structural immunity to partisan manipulation.
